\documentclass[12pt]{report}

 % Add this line to your preamble
% Language setting
% Replace `English' with e.g. `Spanish' to change the document language
\usepackage[english]{babel}

% Set page size and margins
% Replace `letterpaper' with `a4paper' for UK/EU standard size

\usepackage[top=3cm,bottom=3cm,left=3cm,right=3cm,marginparwidth=1.75cm]{geometry}
	
% Load the setspace package
\usepackage{setspace}
\usepackage{amsmath}
\usepackage{listings}
\usepackage{xcolor}
\usepackage{multirow}
\usepackage{amssymb}
\usepackage{pgfplots}
\usepackage{pgfplotstable}
\usepackage{float}
\usepackage{filecontents}
\usepackage{pgfplotstable}

\usepgfplotslibrary{fillbetween}
\usetikzlibrary{arrows.meta, bending}
\pgfplotsset{compat=newest}
\usetikzlibrary{decorations.markings}
\lstdefinestyle{mystyle}{
    commentstyle=\color{gray},
    keywordstyle=\color{blue},
    numberstyle=\tiny\color{gray},
    stringstyle=\color{purple},
    basicstyle=\ttfamily\footnotesize,
    breakatwhitespace=false,         
    breaklines=true,                 
    captionpos=b,                    
    keepspaces=true,                 
    numbers=left,                    
    numbersep=5pt,                  
    showspaces=false,                
    showstringspaces=false,
    showtabs=false,                  
    tabsize=2
}

\lstset{style=mystyle}

% Useful packages
\usepackage{graphicx}
\usepackage[colorlinks=true, allcolors=black]{hyperref}
\usepackage{tikz}
\usepackage{natbib}
\pgfplotsset{width=12cm,compat=1.9}
\setstretch{1.5}

\bibliographystyle{unsrtnat}

\usetikzlibrary{shapes.arrows, positioning, arrows.meta}


\title{Financial Frictions and the Cleansing Effects of Recessions}
\author{Riccardo Dal Cero}


\begin{document}
\section{Free debt case}

\begin{align*}
    \begin{cases}
        V(k_0) = U(d_0^*) + \beta V(k_1), \\
        k_1 = \left[ f(k_{0}) + (1 - \delta) k_{0}  - d^*_0 \right] , \\
        d^*_0 = [\beta V'(k_{1})]^{-1},\\
        k_0 \text{ given.}
    \end{cases}
\end{align*}
The welfare function instead of being log-linear as in the thesis I employ the following quadratic preferences:

\begin{equation}
   W_0= \sum_{t=0}^{\infty}\beta^t\left(\epsilon + \gamma d_t - \frac{\eta}{2}d_t^2 \right)
\end{equation}

where \(\gamma\) and \(\eta\) are assumed to be positive. Moreover in order to be solvable is necessary to approximate
the intertemporal constraint with a linear quadratic approximation.
The steady-state equations are:

\begin{align}
    \frac{1}{\beta}=\alpha Z \hat{k}^{\alpha -1}+1-\delta \label{eq1}\\ 
    \widehat{d} =Z \hat{k}^{\alpha}-\delta \hat{k}\label{eqds}
\end{align}

Using the Taylor theorem to the log utility function:
\begin{equation}
    \ln(d_t) \approx \ln(\hat{d}) + \frac{1}{\hat{d}}(d_t - \hat{d}) - \frac{1}{2\hat{d}^2}(d_t - \hat{d})^2
\end{equation}


Doing the same for the flow of funds constraint:

\begin{align*}
    k_{t+1}= Z \hat{k}^{\alpha}+(1-\delta)\hat{k}+\alpha Z \hat{k}^{\alpha-1}(k_t-\hat{k})+(1-\delta)(k_t-\hat{k})-\hat{d}-(d_t-\hat{d}),\\
    k_{t+1}= Z \hat{k}^{\alpha} - \hat{d}+\left(\alpha Z \hat{k}^{\alpha-1} +1 -\delta\right)(k_t-\hat{k})-(d_t-\hat{d}),
\end{align*}

Using the equation (\ref{eq1}, \ref{eqds}), we get the new :

\begin{align*}
    k_t-\hat{k}=\frac{1}{\beta}(k_t-\hat{k})-(d_t-\hat{d})
\end{align*}

the firm problem in the absence of debt:
\begin{align*}
    W_0= \left(\epsilon + \gamma d_t - \frac{\eta}{2}d_t^2 + \beta V(k_1)\right),
    s.t. k_1 = (1+r)k_0-d_0
\end{align*}
where k and d are deviation from the mean and \(\left(1+r\right)=\frac{1}{\beta}\)

The first order condition is: 
\begin{equation}
    \lambda - \eta d_0^* = \beta V^{\prime}(k_1)
\end{equation}
from the dynamic constraint we obtain \(d_0=(1+r)k_0-k_1\), substituting into the belman equation, we get:
\begin{align*}
    \begin{cases}
        V_(k_0)= \epsilon + \gamma [(1+r)k_0 -k_1] - \frac{\eta}{2}\left[\left(1+r\right)k_0 -k_1\right]^2 +\beta V(k_1)\\
        \gamma -\eta\left[\left(1+r\right)k_0-k_1\right]=\beta V^{\prime}(k_1), \\ 
        k_0 given
    \end{cases}
\end{align*}
Our guess is similar to the quadratic utility function 
\begin{equation}
    V(k_t)=g+h k_t = \frac{m}{2}k_t^2
\end{equation}
Exploiting the guess and the second equation of the above system we get:
\begin{equation}
    k_0=\frac{1}{1+r}\left[\frac{\gamma-\beta h}{\eta}+\frac{\eta-\beta m}{\eta} k_1\right] .
\end{equation}
Substituting the guess into \(V(k_0), V(k_1)\), we get:
\begin{equation}
    \left\{\begin{array}{c}
    g+\frac{h}{1+r}\left(\frac{\gamma-\beta h}{\eta}\right)+\frac{m}{2(1+r)^2}\left(\frac{\gamma-\beta h}{\eta}\right)^2=\varepsilon+\frac{\gamma}{\eta}(\gamma-\beta h)-\frac{(\gamma-\beta h)^2}{2 \eta}+\beta g \\
    \frac{h}{1+r}\left(\frac{\eta-\beta m}{\eta}\right)+\frac{m}{(1+r)^2}\left(\frac{\eta-\beta m}{\eta}\right)\left(\frac{\gamma-\beta h}{\eta}\right)=-\frac{\gamma \beta m}{\eta}+\frac{(\gamma-\beta h) \beta m}{\eta}+\beta h \\
    \frac{m}{2(1+r)^2}\left(\frac{\eta-\beta m}{\eta}\right)^2=-\frac{(\beta m)^2}{2 \eta}+\frac{\beta m}{2}
    \end{array}\right.
\end{equation}
the solutions of the system are:

\begin{align*}
    \begin{cases}
        m=\frac{\eta}{\beta}\left[1-\beta(1+r)^2\right] \\
        h=\frac{\gamma}{\beta} \frac{\left[\beta(1+r)^2-1\right]}{r} \\ 
        g=\frac{\epsilon}{1-\beta}+\frac{\gamma^2}{2 \eta \beta(1-\beta)}\left[\frac{1-\beta(1+r)}{r}\right]^2 \\
    \end{cases}
\end{align*}
Therefore the optimal path for capital is:
\begin{align*}
    k_1 = \frac{k_0 (1 + r) - \frac{\gamma \left( \frac{r - \beta (1 + r)^2 + 1}{r} \right)}{\eta}}{\beta (1 + r)^2}
\end{align*}
Given the substitution \( 1 + r = \frac{1}{\beta} \), we can further simplify the equation for \( k_1 \).

1. Substitute \( 1 + r = \frac{1}{\beta} \) into the expressions for \( m \) and \( h \):

\[ m = \frac{\eta}{\beta} \left[ 1 - \beta \left(\frac{1}{\beta}\right)^2 \right] = \frac{\eta}{\beta} \left[ 1 -
\frac{1}{\beta^2} \right] = \frac{\eta}{\beta} \left[ \frac{\beta^2 - 1}{\beta^2} \right] = \frac{\eta (\beta^2 -
1)}{\beta^3} \]


\[ h = \frac{\gamma}{\beta} \frac{\left[ \beta \left(\frac{1}{\beta}\right)^2 - 1 \right]}{r} = \frac{\gamma}{\beta}
\frac{\left[ \frac{1}{\beta} - 1 \right]}{r} = \frac{\gamma}{\beta} \frac{\left[ 1 - \beta \right]}{r \beta} =
\frac{\gamma (1 - \beta)}{r \beta^2} \]


2. Substitute these values into the equation for \( k_1 \):

The original equation is:

\[ k_1 = \frac{k_0 (1 + r) - \frac{\gamma - \beta h}{\eta}}{1 - \frac{\beta m}{\eta}} \]

Substituting the new values of \( m \) and \( h \):

\[ \gamma - \beta h = \gamma - \beta \left( \frac{\gamma (1 - \beta)}{r \beta^2} \right) = \gamma - \frac{\gamma (1 -
\beta)}{r \beta} = \gamma \left( 1 - \frac{1 - \beta}{r \beta} \right) \]


Given \( r = \frac{1}{\beta} - 1 \):

\[ \gamma \left( 1 - \frac{1 - \beta}{r \beta} \right) = \gamma \left( 1 - \frac{1 - \beta}{(\frac{1}{\beta} - 1) \beta} \right) \]

\[ = \gamma \left( 1 - \frac{1 - \beta}{\frac{1 - \beta \beta}{\beta}} \right) = \gamma \left( 1 - \frac{1 - \beta}{1 -
\beta} \right) = \gamma \left( 1 - 1 \right) = 0 \]


Thus, we see that \( \gamma - \beta h = 0 \).

Now substituting \( m \) and simplifying the denominator:

\[ \frac{\beta m}{\eta} = \frac{\beta \left( \frac{\eta (\beta^2 - 1)}{\beta^3} \right)}{\eta} = \frac{\beta (\beta^2 -
1)}{\beta^3} = \frac{\beta^2 - 1}{\beta^2} \]


\[ 1 - \frac{\beta m}{\eta} = 1 - \frac{\beta^2 - 1}{\beta^2} = \frac{\beta^2 - (\beta^2 - 1)}{\beta^2} =
\frac{1}{\beta^2} \]


Therefore:

\[ k_1 = \frac{k_0 (1 + r) - 0}{\frac{1}{\beta^2}} = k_0 (1 + r) \beta^2 \]

Since \( 1 + r = \frac{1}{\beta} \):

\[ k^{\star}_1 = k_0 \left( \frac{1}{\beta} \right) \beta^2 = \beta k_0 \]

This is the optimal path for capital.
The policy function is:
\[d^{\star}_0=\frac{1}{\beta} k_0 - k_0 \beta = \frac{1 - \beta^2}{\beta} k_0\]
Reminding that \(\beta\) in the steady state of dividends is equal to \((\alpha Z \hat{k}^{\alpha -1}+1-\delta)^{-1}\),
therefore we can rewrite the optimal path of capital as: 
\[k_1 = \beta k_0 = \frac{1}{\alpha Z \hat{k}^{\alpha -1}+1-\delta} k_0 \]
Thus a deviation from the steady state of capital \(k_0\) either positive or negative will follow a converging path to
the steady state since \(\beta\) is less than 1.
The optimal path of dividends is:
\begin{equation}
    d^{\star} = \frac{1-\left(\alpha Z \hat{k}^{\alpha-1}+1-\delta\right)^{-2}}{\left(\alpha Z \hat{k}^{\alpha-1}+1-\delta\right)^{-1}} k_0
\end{equation}
This result suggests that a positive shock to capital will lead to a positive shock to dividends, but the magnitude of
the shock will be smaller than the shock to capital, since \(0<(1-\beta^2)/\beta<1\).
\section{simulation}
Let's use the following parameters:
\begin{table}[H]
    \centering
    \begin{tabular}{lcc}
    \hline Parameter & Symbol & Value \\
    \hline \hline
    Depreciation rate & $\delta$ & 0.07 \\
    Returns to scale & $\alpha$ & 0.70 \\
    Firm productivity & $Z$ & 1 \\
    Discount factor & $\beta$ & 0.95 \\
    \hline
    \end{tabular}
    \caption{Benchmark calibration}
\end{table}

The following plot shows the evolution of capital and dividends over 100 periods, with log deviations from the steady state.

\pgfplotstableread[col sep=comma]{output_data/quadratic_function.csv}\datatable

\begin{figure}[h!]
    \centering
    \begin{tikzpicture}
        \begin{axis}[
            width=12cm,
            height=8cm,
            grid=major,
            xlabel={Periods},
            ylabel={Log Deviations from Steady State},
            legend pos=north east,
            title={Evolution of Capital and Dividends Over Time},
            legend style={font=\small},
            ymajorgrids=true,
            xmajorgrids=true,
        ]
        \addplot[
            color=blue,
        ] table[x=Period,y={Capital (log deviations)}] {\datatable};
        \addlegendentry{Capital (log deviations)}

        \addplot[
            color=green,
        ] table[x=Period,y={Dividends (log deviations)}] {\datatable};
        \addlegendentry{Dividends (log deviations)}
        \end{axis}
    \end{tikzpicture}
    \caption{Log deviations from steady state for capital and dividends over 100 periods.}
    \label{fig:impulse_response1}
\end{figure}
 
The impulse response function \ref{fig:impulse_response1} represents the effect of a positive $10\%$ shock to the steady state level of capital (blue)
and on dividends (green). The rate of convergence is a function of \(\beta\). Lets for example consider the same
positive shock on capital but with a discount factor of 0.9:
\pgfplotstableread[col sep=comma]{output_data/quadratic_function_lowerb.csv}\datatable

\begin{figure}[!h]
    \centering
    \begin{tikzpicture}
        \begin{axis}[
            width=12cm,
            height=8cm,
            grid=major,
            xlabel={Periods},
            ylabel={Log Deviations from Steady State},
            legend pos=north east,
            title={Evolution of Capital and Dividends Over Time},
            legend style={font=\small},
            ymajorgrids=true,
            xmajorgrids=true,
        ]
        \addplot[
            color=blue,
        ] table[x=Period,y={Capital (log deviations)}] {\datatable};
        \addlegendentry{Capital (log deviations)}

        \addplot[
            color=green,
        ] table[x=Period,y={Dividends (log deviations)}] {\datatable};
        \addlegendentry{Dividends (log deviations)}
        \end{axis}
    \end{tikzpicture}
    \caption{Log deviations from steady state for capital and dividends over 100 periods with a discount factor of 0.9.}
    \label{fig:impulse_response2}
\end{figure}
As the impulse response function \ref{fig:impulse_response2} shows, the rate of convergence is faster with a lower discount rate. Therefore,
an increase in the risk free rate will lead to a faster convergence to the steady state level of capital and dividends,
after a positive shock to capital.
Regarding productivity shocks, they affected only the steady state level of capital and dividends, but not the rate of convergence.

\section{Leverage Case}
In this section, I will introduce the possibility for firms to finance their investment through debt. A financial market is
introduced to supply debt to firms, which is assumed to be in a perfectly competitive equilibrium. In this market,
financial intermediaries are risk neutral and compete to lend to firms, while facing an opportunity cost of lending to the risk-free rate.
Moreover, lenders in case of default of the firm which occurs with probability \(1-p\) can seize the firm's output after
facing a cost \(\mu\) for the seizure, which is assumed to be a fraction of the firm's output. Therefore the
participation constraint for the financial intermediaries is formulated as follows:
\begin{equation}
    p R_t B_t + (1-p) \mu Z K_t^{\alpha} = R_f B_t \label{eq:participation}
\end{equation}
where \(R_f\) is the  gross risk-free rate, \(R\) is the gross return on debt, \(B_t\) is the amount of debt, and
\(K_t\) is the amount of capital. Exploiting \ref{eq:partecipation} we can make explicit the gross return on debt:
\begin{equation}
    R_t = \frac{R_f}{p} + \frac{(1-p)\mu Z K_t^{\alpha}}{p B_t}
\end{equation}
The new f-of-f constraint, including the possibility to borrow, is:
\begin{equation}
    K_{t+1} + R_t B_{t} + D_{t} = Z K_t^{\alpha} + B_{t+1} + K_{t} (1-\delta) \label{eq:flow_of_funds}
\end{equation}
As described by the f-of-f constraint, the firm uses the inflows from the output, the debt, and the capital net of
depreciation to finance the capital investment and to pay dividends and debt (interest and principal).
The firm problem is to maximize the discounted utility function over an infinite horizon subject to the budget constraint \ref{eq:flow_of_funds}:
\begin{equation}
    \mathcal{L} = \sum_{t=0}^{\infty} \beta^t \left[ \ln(d_t) + \lambda_t \left( Z K_t^{\alpha} + B_{t+1} - K_{t+1} - \left(\frac{R_f}{p} + \frac{(1-p)\mu Z K_t^{\alpha}}{p B_t}\right) B_t - D_t \right) \right]
\end{equation}
The Euler equations are:
\begin{itemize}
    \item For Dividends:
    \begin{equation}
        \frac{\partial \mathcal{L}}{\partial D_t} = \frac{1}{D_t} - \lambda_t = 0 \implies \lambda_t = \frac{1}{D_t}
    \end{equation}
    \item For Capital:
    \begin{align}
        &\frac{\partial \mathcal{L}}{\partial K_{t+1}} = -\lambda_t + \beta \lambda_{t+1} \alpha Z K_{t+1}^{\alpha-1} 
        \left(\frac{p+(1-p)\mu}{p}\right) = 0 \\
        & \implies \lambda_t = \beta \lambda_{t+1} \alpha Z K_{t+1}^{\alpha-1} \left(\frac{p+(1-p)\mu}{p}\right)
    \end{align}
    \item For Debt:
    \begin{equation}
        \frac{\partial \mathcal{L}}{\partial B_{t+1}} = \lambda_t - \beta \lambda_{t+1} \frac{R_f}{p}= 0 \implies \lambda_t = \beta \lambda_{t+1} \frac{R_f}{p}
    \end{equation}
\end{itemize}
\appendix
\subsection{B: Resolution with Log Utility and Policy Function}

  Consider the utility function \( U(d_t) = \ln(d_t) \). The Bellman equation is:
  
  \[
  V(k_0) = \ln(d_0^*) + \beta V(k_1),
  \]
  
  where \( k_1 = a \lambda k_0 - \lambda d_0^* \) and \( k_0 \) is given.
  
  Assume the value function \( V(k) \) takes the form:
  
  \[
  V(k) = A \ln(k) + B,
  \]
  
  where \( A \) and \( B \) are constants to be determined.
  
  Substitute the hypothesized value function into the Bellman equation:
  
  \[
  A \ln(k_0) + B = \ln(d_0^*) + \beta (A \ln(a \lambda k_0) + \beta B).
  \]
  
  Using \( \ln(a \lambda k_0) = \ln(a \lambda) + \ln(k_0) \):
  
  \[
  A \ln(k_0) + B = \ln(d_0^*) + \beta A \ln(a \lambda) + \beta A \ln(k_0) + \beta B.
  \]
  
  Matching coefficients gives:
  
  \[
  A (1 - \beta) = \beta A \ln(a \lambda),
  \]
  
  \[
  B (1 - \beta) = \ln(d_0^*).
  \]
  
  Solving for \( A \) and \( B \):
  
  \[
  a \lambda = e^{\frac{1 - \beta}{\beta}},
  \]
  
  \[
  B = \frac{\ln(d_0^*)}{1 - \beta}.
  \]
  
  The value function \( V(k_0) \) is:
  
  \[
  V(k_0) = A \ln(k_0) + B,
  \]
  
  where:
  
  \[
  A = \frac{1 - \beta}{\beta \ln(a \lambda)},
  \]
  
  \[
  B = \frac{\ln(d_0^*)}{1 - \beta}.
  \]
  
  The optimal consumption \( d_0^* \) maximizes:
  
  \[
  \ln(d_0^*) + \beta V(a \lambda k_0 - \lambda d_0^*).
  \]
  
  The first-order condition gives:
  
  \[
  d_0^* = \frac{a \lambda k_0}{1 + \beta \lambda}.
  \]
  
  The optimal path for capital is:
  
  \[
  k_{t+1} = a \lambda k_t - \lambda d_t^*,
  \]
  
  substituting \( d_t^* \):
  
  \[
  k_{t+1} = k_t \left( \frac{a \lambda (1 + (\beta - 1) \lambda)}{1 + \beta \lambda} \right).
  \]
  
  This provides the policy function for optimal consumption \( d_0^* \) and the optimal evolution path for capital \( k_t \).
  
\section{Steady-State Computation and Log-Linearization}
\label{sec:steady-state-computation}

This appendix provides the derivation of the steady-state values for capital, dividends, and debt, and the log-linearization of the flow of funds constraint using the first-order Taylor expansion.

\subsection{Steady-State Computation}
Notation marks:
\begin{itemize}
    \item lower case letter represents log deviation from the steady state
    \item upper case letter represents a level variable
    \item \(X^*\) represents the optimal path
\end{itemize}
The utility function is logarithmic:
\begin{equation}
U(D) = \ln(D)
\end{equation}

The flow of funds constraint is:
\begin{equation}
K^{t+1} + \left( \frac{R_f}{p} - \frac{(1-p)}{p} \frac{\mu Z K_t^\alpha}{B} \right) B_t + D_t = Z K^\alpha + B^{t+1}
\end{equation}

At the steady state, we have:
\begin{equation}
K + \left( \frac{R_f}{p} - \frac{(1-p)}{p} \frac{\mu Z K^\alpha}{B} \right) B + D = Z K^\alpha + B
\end{equation}

The Lagrangian for the firm's problem is:
\begin{equation}
\mathcal{L} = \sum_{t=0}^\infty \beta^t \left[ \ln(D_t) - \lambda_t \left( - Z K_t^\alpha - B_{t+1} + K_{t+1} + \left( \frac{R_f}{p} - \frac{(1-p)}{p} \frac{\mu Z K_t^\alpha}{B_t} \right) B_t + D_t \right) \right]
\end{equation}

the first-order conditions are:
\begin{itemize}
    \item For Dividends:
    \begin{equation}
        \frac{\partial \mathcal{L}}{\partial D_t} = \frac{1}{D_t} - \lambda_t = 0 \implies \lambda_t = \frac{1}{D_t}
    \end{equation}
    \item For Capital:
    \begin{align}
        &\frac{\partial \mathcal{L}}{\partial K_{t+1}} = -\lambda_t + \beta \lambda_{t+1} \alpha Z K_{t+1}^{\alpha-1} 
        \left(\frac{p+(1-p)\mu}{p}\right) = 0 \\
        & \implies \lambda_t = \beta \lambda_{t+1} \alpha Z K_{t+1}^{\alpha-1} \left(\frac{p+(1-p)\mu}{p}\right)
    \end{align}
    \item For Debt:
    \begin{equation}
        \frac{\partial \mathcal{L}}{\partial B_{t+1}} = \lambda_t - \beta \lambda_{t+1} \frac{R_f}{p}= 0 \implies \lambda_t = \beta \lambda_{t+1} \frac{R_f}{p}
    \end{equation}
\end{itemize}

Combining the first-order conditions, we get the Euler equations for capital and debt.

\paragraph{Euler Equation for Capital}
\begin{equation}
\frac{1}{D_t} = \beta \frac{1}{D_{t+1}} \alpha Z K_{t+1}^{\alpha-1} \left(\frac{p+(1-p)\mu}{p}\right) \label{eq:euler-capital}
\end{equation}

\paragraph{Euler Equation for Debt}
\begin{equation}
\frac{1}{D_t} = \beta \frac{1}{D_{t+1}} \frac{R_f}{p} \label{eq:euler-debt}
\end{equation}

At the steady state, \(K_t = K_{t+1} = K\), \(D_t = D_{t+1} = D\), and \(B_t = B_{t+1} = B\).

\subsubsection{Steady-State Euler Equation for Capital}
\begin{equation}
K = \left(\frac{\alpha \beta Z (p+(1-p)\mu)}{p}\right)^{\frac{1}{1-\alpha}}
\end{equation}

\subsubsection{Steady-State Euler Equation for Debt}
\begin{equation}
p = \frac{\beta}{R_f}
\end{equation}


\subsubsection{Steady-State Flow of Funds Constraint}
Using the flow of funds constraint at the steady state:
\begin{equation}
K + \left( \frac{R_f}{p} - \frac{(1-p)}{p} \frac{\mu Z K^\alpha}{B} \right) B + D = Z K^\alpha + B
\end{equation}

we get the steady state level of dividends:
\begin{equation}
D = Z K^\alpha \left(\frac{p + \mu (1-p)}{p}\right) - K - B \left(\frac{R_f - p}{p}\right)
\end{equation}

\subsection{Log-Linearization Using Taylor Expansion}

We denote deviations from the steady state using lower-case letters for the log variables and upper-case letters for the level variables.

\subsubsection{Flow of Funds Constraint}

The original flow of funds constraint is:
\begin{equation}
K_{t+1} + \left( \frac{R_f}{p} - \frac{(1-p)}{p} \frac{\mu F(K_t)}{B_t} \right) B_t + D_t = F(K_t) + B_{t+1}
\end{equation}
where $F(K_t) = Z K_t^\alpha$.

At the steady state, we have:
\begin{equation}
K + \left( \frac{R_f}{p} - \frac{(1-p)}{p} \frac{\mu Z K^\alpha}{B} \right) B + D = Z K^\alpha + B
\end{equation}

Using log deviations:
\begin{equation}
k_t = \ln\left(\frac{K_t}{K}\right), \quad d_t = \ln\left(\frac{D_t}{D}\right), \quad b_t = \ln\left(\frac{B_t}{B}\right), \quad k_{t+1} = \ln\left(\frac{K_{t+1}}{K}\right), \quad b_{t+1} = \ln\left(\frac{B_{t+1}}{B}\right)
\end{equation}

\subsubsection{Linearizing the Flow of Funds Constraint}


Using a first-order Taylor expansion around the steady state (\(K, B, D\)), we get:
\begin{equation}
    K_{t+1} = f(K) + (1-R)B - D + \mathbb{J}^{T}\begin{bmatrix}
        K_t-K \\
        B_t-B \\
        B_{t+1}-B \\
        D_t-D
\end{bmatrix}
\end{equation}

where:
\begin{equation}
    \mathbb{J} = \begin{bmatrix}
        \frac{\partial K_{t+1}}{\partial K_t} \\
        \frac{\partial K_{t+1}}{\partial B_t}\\
        \frac{\partial K_{t+1}}{\partial B_{t+1}}\\
        \frac{\partial K_{t+1}}{\partial D_t}
    \end{bmatrix}
\end{equation}
exploiting the Euler equations for capital and debt \ref{eq:euler-capital} and \ref{eq:euler-debt} we get:
\begin{equation}
    (K_{t+1}-K) - (B_{t+1}-B) = \frac{1}{\beta} \left[\left(K_t - K\right) - \left(B_t - B\right)\right] - \left(D_t - D\right)
\end{equation}
rewriting using lowercase letters to express the log deviation from the steady state:
\begin{equation}
    k_{t+1} - b{t+1} = \frac{1}{\beta} \left(k_t - b_t\right) - d_t
\end{equation}
we can rewrite expressing the equity as (\(n=k-b\)):
\begin{equation}
    n_{t+1} = \frac{1}{\beta} n_t - d_t
\end{equation}
This completes the log-linearization of the flow of funds constraint around the steady state.




 

\end{document}